\documentclass[10pt,a4paper]{amsart}
\usepackage{amsmath,amssymb,amsthm,graphicx}
\usepackage{verbatim}
\newenvironment{reference}{\comment}{\endcomment}
%\newenvironment{reference}{}{}
\newenvironment{slogan}{\comment}{\endcomment}
\newenvironment{history}{\comment}{\endcomment}

% For commutative diagrams we use Xy-pic
\usepackage[all]{xy}

% We use 2cell for 2-commutative diagrams.
\xyoption{2cell}
% Theorem environments.
%
\theoremstyle{plain}
\newtheorem{theorem}[subsection]{Theorem}
\newtheorem{proposition}[subsection]{Proposition}
\newtheorem{lemma}[subsection]{Lemma}

\theoremstyle{definition}
\newtheorem{definition}[subsection]{Definition}
\newtheorem{example}[subsection]{Example}
\newtheorem{exercise}[subsection]{Exercise}
\newtheorem{situation}[subsection]{Situation}

\theoremstyle{remark}
\newtheorem{remark}[subsection]{Remark}
\newtheorem{remarks}[subsection]{Remarks}

\numberwithin{equation}{subsection}

% Macros
%
\def\lim{\mathop{\mathrm{lim}}\nolimits}
\def\colim{\mathop{\mathrm{colim}}\nolimits}
\def\Spec{\mathop{\mathrm{Spec}}}
\def\Hom{\mathop{\mathrm{Hom}}\nolimits}
\def\Ext{\mathop{\mathrm{Ext}}\nolimits}
\def\SheafHom{\mathop{\mathcal{H}\!\mathit{om}}\nolimits}
\def\SheafExt{\mathop{\mathcal{E}\!\mathit{xt}}\nolimits}
\def\Sch{\mathit{Sch}}
\def\Mor{\mathop{\mathrm{Mor}}\nolimits}
\def\Ob{\mathop{\mathrm{Ob}}\nolimits}
\def\Sh{\mathop{\mathit{Sh}}\nolimits}
\def\NL{\mathop{N\!L}\nolimits}
\def\CH{\mathop{\mathrm{CH}}\nolimits}
\def\proetale{{pro\text{-}\acute{e}tale}}
\def\etale{{\acute{e}tale}}
\def\QCoh{\mathit{QCoh}}
\def\Ker{\mathop{\mathrm{Ker}}}
\def\Im{\mathop{\mathrm{Im}}}
\def\Coker{\mathop{\mathrm{Coker}}}
\def\Coim{\mathop{\mathrm{Coim}}}

% Boxtimes
%
\DeclareMathSymbol{\boxtimes}{\mathbin}{AMSa}{"02}

%
% Macros for moduli stacks/spaces
%
\def\QCohstack{\mathcal{QC}\!\mathit{oh}}
\def\Cohstack{\mathcal{C}\!\mathit{oh}}
\def\Spacesstack{\mathcal{S}\!\mathit{paces}}
\def\Quotfunctor{\mathrm{Quot}}
\def\Hilbfunctor{\mathrm{Hilb}}
\def\Curvesstack{\mathcal{C}\!\mathit{urves}}
\def\Polarizedstack{\mathcal{P}\!\mathit{olarized}}
\def\Complexesstack{\mathcal{C}\!\mathit{omplexes}}
% \Pic is the operator that assigns to X its picard group, usage \Pic(X)
% \Picardstack_{X/B} denotes the Picard stack of X over B
% \Picardfunctor_{X/B} denotes the Picard functor of X over B
\def\Pic{\mathop{\mathrm{Pic}}\nolimits}
\def\Picardstack{\mathcal{P}\!\mathit{ic}}
\def\Picardfunctor{\mathrm{Pic}}
\def\Deformationcategory{\mathcal{D}\!\mathit{ef}}
\usepackage{cleveref}
\crefname{originallemma}{original lemma}{original lemmas}
\Crefname{originallemma}{Original lemma}{Original lemmas}
\makeatletter
\newlabel{lemma-pullback-along-p-power-cover}{{07Q9}{source}{}{sourcepointer-1}{}}
\newlabel{lemma-pullback-along-p-power-cover@cref}{{[originallemma][0][]07Q9}{[1][0][]source}{}{}{}}
\makeatother
\crefname{originallemma}{original lemma}{original lemmas}
\Crefname{originallemma}{Original lemma}{Original lemmas}
\makeatletter
\newlabel{lemma-pullback-along-p-power-cover-cohomology}{{07QA}{source}{}{sourcepointer-2}{}}
\newlabel{lemma-pullback-along-p-power-cover-cohomology@cref}{{[originallemma][0][]07QA}{[1][0][]source}{}{}{}}
\makeatother
\crefname{originalreference}{original reference}{original references}
\Crefname{originalreference}{Original reference}{Original references}
\makeatletter
\newlabel{morphisms-lemma-smooth-etale-over-affine-space}{{054L}{source}{}{sourcepointer-3}{}}
\newlabel{morphisms-lemma-smooth-etale-over-affine-space@cref}{{[originalreference][0][]054L}{[1][0][]source}{}{}{}}
\makeatother
\crefname{originalreference}{original reference}{original references}
\Crefname{originalreference}{Original reference}{Original references}
\makeatletter
\newlabel{etale-lemma-relative-frobenius-etale}{{0EBS}{source}{}{sourcepointer-4}{}}
\newlabel{etale-lemma-relative-frobenius-etale@cref}{{[originalreference][0][]0EBS}{[1][0][]source}{}{}{}}
\makeatother
\title{Relative Frobenius is an iterated alpha\_p-cover of degree p\textasciicircum{}d}
\author{The Stacks Project Authors}
\date{stacks-project:crystalline.tex:local-snapshot-20261012; GNU FDL 1.2 or later}
\begin{document}
\maketitle
\noindent\textit{Editorial note.} All delivered mathematical text is the author's own, copied verbatim from the pinned Stacks Project chapter file crystalline.tex; no statement or proof step was rewritten, reordered or removed.\par\medskip
\noindent\textit{Editorial note.} Span recomputation by label pairing (LANE\_BRIEF 10.2): in this chapter the author's proof environment follows the statement environment as a standalone block; the spec covers the complete environment pair (statement span, proof span) as recorded in the pool candidate.\par\medskip
\ifdefined\hypertarget\hypertarget{sourcepointer-1}{}\fi
\noindent\textit{Source reference.} Original lemma 07Q9: sources/stacks-chow/tags.txt, source lines 9824--9824. Cross-chapter standard prerequisite cited by the author inside the counted span; it is used as a pointer and its own proof is not part of this item.\par\medskip
\ifdefined\hypertarget\hypertarget{sourcepointer-2}{}\fi
\noindent\textit{Source reference.} Original lemma 07QA: sources/stacks-chow/tags.txt, source lines 9825--9825. Cross-chapter standard prerequisite cited by the author inside the counted span; it is used as a pointer and its own proof is not part of this item.\par\medskip
\ifdefined\hypertarget\hypertarget{sourcepointer-3}{}\fi
\noindent\textit{Source reference.} Original reference 054L: sources/stacks-chow/tags.txt, source lines 6644--6644. Cross-chapter standard prerequisite cited by the author inside the counted span; it is used as a pointer and its own proof is not part of this item.\par\medskip
\ifdefined\hypertarget\hypertarget{sourcepointer-4}{}\fi
\noindent\textit{Source reference.} Original reference 0EBS: sources/stacks-chow/tags.txt, source lines 17932--17932. Cross-chapter standard prerequisite cited by the author inside the counted span; it is used as a pointer and its own proof is not part of this item.\par\medskip
\begin{lemma}
\label{lemma-pullback-relative-frobenius}
In the situation above, assume that $X \to S_0$ is smooth of relative
dimension $d$. Then $F_{X/S_0}$ is an iterated $\alpha_p$-cover
of degree $p^d$. Hence Lemmas \ref{lemma-pullback-along-p-power-cover} and
\ref{lemma-pullback-along-p-power-cover-cohomology} apply to this
situation. In particular, for any crystal in quasi-coherent modules
$\mathcal{G}$ on $\text{Cris}(X^{(1)}/S)$ the map
$$
F_{X/S_0}^* : H^i(\text{Cris}(X^{(1)}/S), \mathcal{G})
\longrightarrow
H^i(\text{Cris}(X/S), F_{X/S_0, \text{cris}}^*\mathcal{G})
$$
has kernel and cokernel annihilated by $p^{d(i + 1)}$.
\end{lemma}

\begin{proof}
It suffices to prove the first statement. To see this we may assume
that $X$ is \'etale over $\mathbf{A}^d_{S_0}$, see
Morphisms, Lemma \ref{morphisms-lemma-smooth-etale-over-affine-space}.
Denote $\varphi : X \to \mathbf{A}^d_{S_0}$ this \'etale morphism.
In this case the relative Frobenius of $X/S_0$ fits into a diagram
$$
\xymatrix{
X \ar[d] \ar[r] & X^{(1)} \ar[d] \\
\mathbf{A}^d_{S_0} \ar[r] & \mathbf{A}^d_{S_0}
}
$$
where the lower horizontal arrow is the relative frobenius morphism
of $\mathbf{A}^d_{S_0}$ over $S_0$. This is the morphism which raises
all the coordinates to the $p$th power, hence it is an iterated
$\alpha_p$-cover. The proof is finished by observing that the diagram
is a fibre square, see
\'Etale Morphisms, Lemma \ref{etale-lemma-relative-frobenius-etale}.
\end{proof}

\end{document}
